\chapter{Hyperparameter Grids}
\label{app:hyperparameter-grids}

This appendix reports the configurations selected by the validation-only
model-selection searches summarized in \cref{sec:hyperparameter-search}. The
tables list the archived selection score used by each search. Final external
test metrics and switch metrics are reported in \cref{chap:results}; they should
be considered the primary performance results.

For current-level persistence, some duration metrics were refreshed without
retraining after dataset-alignment updates. The configurations below are the
selected model artifacts; the refreshed validation and test metrics used for
the final discussion are those in \cref{chap:results}.

\section{Search Design Summary}

\Cref{tab:appendix-grid-search-summary} summarizes the model-selection
searches used in the thesis. All searches use validation-only criteria. The
external holdout split is not used to choose hyperparameters, thresholds, or
checkpoints.
The purpose of this table is to make the scale of the model-selection phase
explicit before listing the selected configurations. It reports the criterion
used to rank trials, the number of evaluated trials, and the main
hyperparameter families varied by each search.

\begin{table}[htbp]
  \centering
  \caption{Summary of the hyperparameter searches. Trial counts refer to the archived searches used for the reported model artifacts.}
  \label{tab:appendix-grid-search-summary}
  \scriptsize
  \setlength{\tabcolsep}{3pt}
  \renewcommand{\arraystretch}{1.15}
  \begin{tabular}{
    @{}
    >{\raggedright\arraybackslash}p{0.20\textwidth}
    >{\raggedright\arraybackslash}p{0.22\textwidth}
    >{\raggedright\arraybackslash}p{0.18\textwidth}
    >{\raggedleft\arraybackslash}p{0.08\textwidth}
    >{\raggedright\arraybackslash}p{0.23\textwidth}
    @{}
  }
    \toprule
    Task & Search & Selection criterion & Trials & Main varied hyperparameters \\
    \midrule
    Autoregressive
      & GRU S2V and encoder-decoder, raw and context-standard variants
      & validation RMSE, minimized
      & 576
      & hidden size, number of layers, learning rate, weight decay, teacher forcing, architecture. \\
    Autoregressive
      & PatchTST, raw and context-standard variants
      & validation RMSE, minimized
      & 1728
      & patch length, stride, model dimension, attention heads, number of layers, dropout, learning rate. \\
    Autoregressive
      & XGBoost multi-output horizon regressors
      & validation RMSE, minimized
      & 2304
      & number of trees, depth, learning rate, subsampling, column sampling, minimum child weight, \(L_2\) regularization. \\
    Current-level persistence
      & Learnable-shapelet MLP, Transformer, and convolution heads
      & validation MAE in seconds, minimized
      & 648
      & shapelet count, hidden dimensions, scalar encoder size, depth, dropout, learning rate. \\
    Current-level persistence
      & XGBoost duration regressor
      & validation MAE in seconds, minimized
      & 1152
      & number of trees, depth, learning rate, subsampling, column sampling, minimum child weight, \(L_2\) regularization. \\
    Long-fade detection
      & XGBoost, TCN, and shapelet-convolution classifiers
      & validation AUPRC, maximized
      & 8928
      & tree parameters, TCN channels/depth/kernel/dropout, shapelet count/lengths, scalar hidden dimension, learning rate, batch size. \\
    Survival persistence
      & XGBoost-AFT
      & validation AFT negative log-likelihood, minimized
      & 1920
      & input representation, sample weighting, AFT distribution, AFT scale, tree depth, learning rate, subsampling, column sampling, regularization. \\
    Survival persistence
      & Discrete-time TCN
      & validation integrated Brier score, minimized
      & 1620
      & binning strategy, input representation, sample weighting, channel schedule, kernel size, dropout, learning rate, weight decay, batch size. \\
    \bottomrule
  \end{tabular}
\end{table}

For tree-based grids, the configuration files use explicit numeric ranges
defined by start, stop, and step values. For neural grids, categorical lists are
used for architecture choices such as hidden dimensions, layer counts, dropout
values, and learning rates. The resulting trial tables are archived under
\path{results/grid_searches/}; the tables below report only the selected
configuration for each model family or variant.

\section{Autoregressive Forecasting}

The autoregressive searches are run separately for the raw and
context-standardized variants. Chronos is excluded from this table because it
is evaluated in zero-shot mode and no hyperparameters are fitted.
\Cref{tab:appendix-ar-best-configs} reports the best validation-selected
configuration for each trainable autoregressive model and input variant. The
selection score is the validation RMSE on the raw signal scale; the external
test metrics are discussed in the results chapter rather than repeated here.

\begin{table}[htbp]
  \centering
  \caption{Best validation-selected autoregressive configurations.}
  \label{tab:appendix-ar-best-configs}
  \scriptsize
  \setlength{\tabcolsep}{3pt}
  \renewcommand{\arraystretch}{1.18}
  \begin{tabular}{
    @{}
    >{\raggedright\arraybackslash}p{0.17\textwidth}
    >{\raggedright\arraybackslash}p{0.13\textwidth}
    >{\raggedright\arraybackslash}p{0.48\textwidth}
    >{\raggedleft\arraybackslash}p{0.14\textwidth}
    @{}
  }
    \toprule
    Model & Variant & Selected configuration & Validation RMSE \\
    \midrule
    GRU S2V & raw
      & hidden size 16; 1 layer; learning rate 0.001; weight decay 0; teacher forcing 0
      & 5.520 \\
    GRU S2V & context-standard
      & hidden size 64; 1 layer; learning rate 0.001; weight decay 0.0001; teacher forcing 0
      & 3.716 \\
    GRU encoder-decoder & raw
      & hidden size 64; 1 layer; learning rate 0.001; weight decay 0.0001; teacher forcing 0.5
      & 4.023 \\
    GRU encoder-decoder & context-standard
      & hidden size 64; 3 layers; learning rate 0.0003; weight decay 0; teacher forcing 0.25
      & 3.691 \\
    PatchTST & raw
      & patch length 10; stride 2; model dimension 128; 4 heads; 3 layers; dropout 0.1; learning rate 0.001
      & 5.526 \\
    PatchTST & context-standard
      & patch length 5; stride 5; model dimension 32; 2 heads; 1 layer; dropout 0.1; learning rate 0.001
      & 3.815 \\
    XGBoost & raw
      & 300 trees; max depth 2; learning rate 0.03; subsample 1.0; column sample 0.8; minimum child weight 1; \(L_2=2.0\)
      & 6.381 \\
    XGBoost & context-standard
      & 300 trees; max depth 4; learning rate 0.03; subsample 1.0; column sample 0.8; minimum child weight 3; \(L_2=2.0\)
      & 3.912 \\
    \bottomrule
  \end{tabular}
\end{table}

\clearpage

\section{Current-Level Persistence}

\Cref{tab:appendix-clp-best-configs} reports the selected duration-regression
configuration for each current-level persistence model. The target for these
models is the remaining time for which the signal stays above the
current-level threshold \(S_t-\delta\). The table lists validation MAE in
seconds because this was the selection metric used by the duration-regression
searches.

\begin{table}[htbp]
  \centering
  \caption{Best validation-selected current-level persistence configurations.}
  \label{tab:appendix-clp-best-configs}
  \scriptsize
  \setlength{\tabcolsep}{3pt}
  \renewcommand{\arraystretch}{1.18}
  \begin{tabular}{
    @{}
    >{\raggedright\arraybackslash}p{0.22\textwidth}
    >{\raggedright\arraybackslash}p{0.54\textwidth}
    >{\raggedleft\arraybackslash}p{0.16\textwidth}
    @{}
  }
    \toprule
    Model & Selected configuration & Validation MAE [s] \\
    \midrule
    Shapelet MLP
      & hidden dimension 64; 2 hidden layers; dropout 0.1; 64 shapelets per length; scalar encoder dimension 32; learning rate 0.001
      & 5590.139 \\
    Shapelet Transformer
      & model dimension 32; 1 Transformer layer; dropout 0; 32 shapelets per length; scalar encoder dimension 16; learning rate 0.001
      & 5653.660 \\
    Shapelet Convolution
      & convolution channels 64; 2 convolution layers; dropout 0; 16 shapelets per length; scalar encoder dimension 32; learning rate 0.001
      & 5408.826 \\
    XGBoost duration regressor
      & 1200 trees; max depth 6; learning rate 0.11; subsample 0.8; column sample 0.8; minimum child weight 3; \(L_2=1.0\)
      & 5465.665 \\
    \bottomrule
  \end{tabular}
\end{table}

\clearpage

\section{Long-Fade Detection}

The long-fade detection search is a classification search. It compares one
tabular gradient-boosted model and two neural sequence models on the same
event-level binary target. The selection criterion is validation AUPRC because
the task is imbalanced and the positive class corresponds to operationally
important long fade events. The table reports the selected configuration for
each classifier family.
\Cref{tab:appendix-lfd-best-configs} therefore should be read as a summary of
the best event-level classifiers, not as a switch-performance table. Switch
quality is computed only after converting classifier probabilities into
post-processed switch sequences.

\begin{table}[htbp]
  \centering
  \caption{Best validation-selected long-fade detection configurations.}
  \label{tab:appendix-lfd-best-configs}
  \scriptsize
  \setlength{\tabcolsep}{3pt}
  \renewcommand{\arraystretch}{1.18}
  \begin{tabular}{
    @{}
    >{\raggedright\arraybackslash}p{0.24\textwidth}
    >{\raggedright\arraybackslash}p{0.52\textwidth}
    >{\raggedleft\arraybackslash}p{0.16\textwidth}
    @{}
  }
    \toprule
    Model & Selected configuration & Validation AUPRC \\
    \midrule
    XGBoost classifier
      & 200 trees; max depth 6; learning rate 0.1; subsample 0.8; column sample 0.9; \(L_2=0.1\)
      & 0.999816 \\
    TCN classifier
      & 128 hidden channels; 5 temporal blocks; kernel size 3; dropout 0; scalar hidden dimension 64; learning rate 0.0003; batch size 128
      & 0.999972 \\
    Shapelet convolution classifier
      & shapelet lengths 10, 15, and 20; 48 shapelets per length; 32 convolution channels; 1 convolution layer; dropout 0; scalar hidden dimension 32; learning rate 0.003
      & 0.999949 \\
    \bottomrule
  \end{tabular}
\end{table}

\clearpage

\section{Survival Persistence}

The survival-persistence searches are reported separately from the
classification and regression tasks because their validation criteria are
survival-specific. XGBoost-AFT is selected by validation AFT negative
log-likelihood on continuous-time censored targets. The discrete-time TCN is
selected by validation integrated Brier score after converting the same
continuous-time labels into discrete survival bins during training.
\Cref{tab:appendix-survival-best-configs} lists the selected survival models
and their task-native validation criterion. The two scores are not directly
comparable because XGBoost-AFT and the discrete-time TCN optimize different
survival objectives.

\begin{table}[htbp]
  \centering
  \caption{Best validation-selected survival persistence configurations.}
  \label{tab:appendix-survival-best-configs}
  \scriptsize
  \setlength{\tabcolsep}{3pt}
  \renewcommand{\arraystretch}{1.18}
  \begin{tabular}{
    @{}
    >{\raggedright\arraybackslash}p{0.22\textwidth}
    >{\raggedright\arraybackslash}p{0.49\textwidth}
    >{\raggedright\arraybackslash}p{0.15\textwidth}
    >{\raggedleft\arraybackslash}p{0.09\textwidth}
    @{}
  }
    \toprule
    Model & Selected configuration & Selection metric & Score \\
    \midrule
    XGBoost-AFT
      & raw flattened context; uniform sample weights; logistic AFT loss with scale 1.0; 1200 boosting rounds; max depth 3; learning rate 0.1; \(L_2=5.0\)
      & validation AFT negative log-likelihood
      & 1.610 \\
    Discrete-time TCN
      & relative-to-current sequence with scalar context; hybrid survival bins; channels 32, 64, and 64; kernel size 3; dilations 1, 2, 4, 8, and 16; dropout 0.25; learning rate 0.0015
      & validation Integrated Brier Score
      & 0.104 \\
    \bottomrule
  \end{tabular}
\end{table}
